% 阿诺德·索末菲（综述）
% license CCBYSA3
% type Wiki

本文根据 CC-BY-SA 协议转载翻译自维基百科\href{https://en.wikipedia.org/wiki/Arnold_Sommerfeld}{相关文章}

阿诺德·约翰内斯·威廉·索末菲（Arnold Johannes Wilhelm Sommerfeld，德语发音：[ˈaʁnɔlt joˈhanəs ˈvɪlhɛlm ˈzɔmɐˌfɛlt]；1868年12月5日－1951年4月26日）是一位德国理论物理学家，他在原子物理与量子物理的发展中开创了先河，同时还培养并指导了新一代理论物理学的重要学生。他曾担任7位诺贝尔奖得主的博士导师与博士后导师，并指导过至少30位其他著名的物理学家和化学家。只有J. J. 汤姆孙的学生群体才能在成就上与之相提并论。

他引入了第二个量子数——方位量子数，以及第三个量子数——磁量子数。他还提出了精细结构常数，并开创了X射线波动理论。
\subsection{早年生活与教育}
索末菲尔德于1868年出生在一个在普鲁士有着深厚家族根基的家庭。他的母亲策齐莉·马蒂亚斯（Cäcilie Matthias，1839–1902）是波茨坦一位建筑商的女儿。他的父亲弗朗茨·索末菲尔德（Franz Sommerfeld，1820–1906）是一名医生，出身于柯尼斯堡的一个显赫家庭。1822年，阿诺德的祖父自内陆迁居到柯尼斯堡，进入普鲁士王国的宫廷邮政秘书处任职。索末菲尔德在其家族所属的普鲁士福音派新教教会中接受基督教洗礼。虽然他本人并不虔信宗教，但他终生未曾放弃自己的基督信仰。\(^\text{[3][4][5]}\)

索末菲尔德在故乡东普鲁士柯尼斯堡的阿尔贝蒂纳大学学习数学和自然科学。他的博士论文导师是数学家费迪南德·冯·林德曼，\(^\text{[5]}\)同时他还受益于数学家阿道夫·胡尔维茨和大卫·希尔伯特以及物理学家埃米尔·维歇特的授课。\(^\text{[7]}\)他曾参加学生兄弟会“德意志大学生社团”，并因此在决斗中在脸上留下了一道伤疤。\(^\text{[8]}\)1891年10月24日（时年22岁），他获得了博士学位。\(^\text{[9]}\)

获得博士学位后，索末菲尔德留在柯尼斯堡攻读教师资格。他于1892年通过了全国考试，随后开始了一年的兵役，服役于柯尼斯堡的预备役部队。他于1893年9月完成义务兵役，之后八年中，每年继续自愿服役八周。凭借翘起的胡须、健壮的体格、典型的普鲁士军人风度，以及脸上的剑伤，他给人留下了如同骠骑兵上校般的印象。\(^\text{[8]}\)
\subsection{职业生涯}
\subsubsection{哥廷根}
1893年10月，索末菲尔德前往哥廷根大学，该校当时是德国数学研究的中心。\(^\text{[10]}\)在那里，他通过一次幸运的私人关系，成为矿物学研究所西奥多·李比施（Theodor Liebisch）的助理——李比施曾是柯尼斯堡大学的教授，也是索末菲尔德家族的朋友。\(^\text{[11]}\)

1894年9月，索末菲尔德成为费利克斯·克莱因（Felix Klein）的助理，负责在克莱因的讲座期间进行全面的笔记整理，并为数学阅览室撰写讲义，同时还管理阅览室。\(^\text{[8]}\)在克莱因的指导下，索末菲尔德于1895年完成了他的《教授资格论文》，这使他得以在哥廷根大学成为一名私人讲师。\(^\text{[13]}\)

作为一名私人讲师，索末菲尔德讲授了广泛的数学及数理物理课程。他最早在哥廷根开设了偏微分方程课程，\(^\text{[8]}\)并在其整个教学生涯中不断发展，最终成为其教材系列《理论物理讲义》的第六卷，题为《物理中的偏微分方程》。\(^\text{[14]}\)
\begin{figure}[ht]
\centering
\includegraphics[width=6cm]{./figures/6f359ea6e05a3107.png}
\caption{1897年的索末菲尔德} \label{fig_ANOSf_1}
\end{figure}
1895年与1896年，克莱因关于旋转体的讲座促使克莱因与索末菲尔德合作撰写了四卷本著作《陀螺理论》——这一合作历时13年（1897–1910）。前两卷探讨理论，后两卷则涉及其在地球物理学、天文学和技术中的应用。\(^\text{[8]}\)索末菲尔德与克莱因的合作，使他在思维方式上更倾向于应用数学，并在讲授艺术上深受影响。\(^\text{[14]}\)

在哥廷根期间，索末菲尔德结识了约翰娜·黑普夫纳），她是哥廷根博物馆馆长\(^\text{[15]}\)恩斯特·黑普夫纳的女儿。1897年10月，索末菲尔德开始担任克劳斯塔尔-策勒费尔德矿业学院的数学教席，接替威廉·维恩。这一职位的收入足以使他最终与约翰娜结婚。\(^\text{[7][8][10]}\)

应克莱因的请求，索末菲尔德还担任了《数学科学百科全书》第五卷的主编。这是一项重大的工作，从1898年持续到1926年。\(^\text{[8][13]}\)
\subsubsection{亚琛}
1900年，在克莱因的推动下，索末菲尔德受聘为皇家亚琛工业大学（Königliche Technische Hochschule Aachen，后来的亚琛工业大学 RWTH Aachen University）应用力学教席的特别教授。在亚琛期间，他发展了流体力学理论，并对此保持了长期的研究兴趣。后来，在慕尼黑大学，索末菲尔德的学生路德维希·霍普夫和维尔纳·海森堡分别以这一课题撰写了他们的博士论文。\(^\text{[7][8][10][13]}\)

在亚琛任职期间，索末菲尔德对滑动轴承润滑理论的研究作出了重要贡献，因此被邓肯·道森列为23位“摩擦学先驱人物”之一。\(^\text{[16]}\)
\subsubsection{慕尼黑}
自1906年起，索末菲尔德出任慕尼黑大学物理学正教授，并担任新成立的理论物理研究所所长。这一职位是由时任慕尼黑物理研究所所长的威廉·伦琴（Wilhelm Röntgen）为他推荐的。\(^\text{[17]}\)在索末菲尔德看来，这一任命仿佛是被召唤进入一个“特权的行动领域”。\(^\text{[14]}\)

直到19世纪末、20世纪初，实验物理在德国学界中一直被认为地位更高。然而，20世纪初，慕尼黑的索末菲尔德与哥廷根大学的马克斯·玻恩等理论家，凭借其早期的数学训练，扭转了这种局面，使数理物理，即理论物理，成为推动学科发展的主导力量，而实验物理则用于验证或推动理论的进展。\(^\text{[18]}\)

在索末菲尔德门下获得博士学位后，沃尔夫冈·泡利、维尔纳·海森堡和瓦尔特·海特勒先后成为玻恩的助手，并在当时快速发展的量子力学领域作出了重要贡献。

在慕尼黑长达32年的教学生涯中，索末菲尔德既讲授普通课程和专业课程，也主持研讨班与学术讨论会。\(^\text{[14]}\)普通课程内容包括：力学、可变形体力学、电动力学、光学、热力学与统计力学，以及物理中的偏微分方程。这些课程每周上课四小时，冬季学期13周，夏季学期11周，面向已经修读过伦琴（后来是威廉·维恩）所开设的实验物理课程的学生。此外，每周还有两个小时的课堂专门用于问题讨论。专业课程则紧扣研究热点，往往与索末菲尔德的研究兴趣相关；这些课程的内容后来常常以科学论文的形式发表。专业讲座的目标在于直面理论物理中的前沿问题，使索末菲尔德和学生能够对问题形成系统理解，而不论他们是否能真正解决这些前沿提出的难题。\(^\text{[19]}\)在研讨班与学术讨论会期间，学生会被指定阅读当下的最新文献，并准备口头报告。\(^\text{[14]}\)自1942年至1951年，索末菲尔德致力于整理自己的讲义，以便出版。\(^\text{[18]}\)这些讲义最终以六卷本《理论物理讲义》的形式问世。

关于学生名单，请参见按类别整理的列表。\(^\text{[20]}\)索末菲尔德的四位博士生——维尔纳·海森堡、沃尔夫冈·泡利、彼得·德拜和汉斯·贝特后来都获得了诺贝尔奖。而其他学生中，最为著名的包括瓦尔特·海特勒、鲁道夫·佩耶尔斯、\(^\text{[22]}\)卡尔·贝歇特、赫尔曼·布吕克、保罗·彼得·埃瓦尔德、尤金·芬伯格、\(^\text{[23]}\)赫伯特·弗勒利希、欧文·福厄斯、恩斯特·吉耶尔曼、赫尔穆特·黑内尔、路德维希·霍普夫、阿道夫·克拉策尔、奥托·拉波尔特、威廉·伦茨、卡尔·迈斯纳、\(^\text{[24]}\)鲁道夫·泽利格尔、恩斯特·施蒂克尔伯格、海因里希·韦尔克、格雷戈尔·文策尔、阿尔弗雷德·朗德以及莱昂·布里渊，\(^\text{[25]}\)他们都各自成名。在索末菲尔德指导下的三位博士后——莱纳斯·鲍林、\(^\text{[26]}\)伊西多·拉比\(^\text{[27]}\)和马克斯·冯·劳厄\(^\text{[28]}\)——也获得了诺贝尔奖。此外，还有十位博士后，包括威廉·阿利斯、\(^\text{[29]}\)爱德华·康登、\(^\text{[30]}\)卡尔·埃卡特、\(^\text{[31]}\)埃德温·C·肯布尔、\(^\text{[32]}\)威廉·V·休斯顿、\(^\text{[33]}\)卡尔·赫茨菲尔德、\(^\text{[34]}\)瓦尔特·科塞尔、菲利普·M·莫尔斯、\(^\text{[35][36]}\) 霍华德·罗伯逊\(^\text{[37]}\) 以及沃伊切赫·鲁比诺维奇，\(^\text{[38]}\)他们同样在学术界各自取得了成就。索末菲尔德在亚琛工业大学的本科生瓦尔特·罗戈夫斯基后来也声名显赫。马克斯·玻恩曾认为，索末菲尔德的能力之一在于“发现和培养人才”。\(^\text{[39]}\)阿尔伯特·爱因斯坦对他说过：“我特别钦佩你的一点是，你似乎能从土地里敲打出这么多年轻的人才。”\(^\text{[39]}\)作为教授与研究所所长，索末菲尔德的风格并没有在他与同事和学生之间制造隔阂。他鼓励他们的合作，而他们的想法往往也影响了他本人的物理观念。他经常在家中款待他们，并在研讨班和学术讨论会前后与他们在咖啡馆聚会。索末菲尔德还拥有一间高山滑雪小屋，他常邀请学生们前往，在那里对物理的讨论往往和这项运动一样充满挑战。\(^\text{[40]}\)

在慕尼黑期间，索末菲尔德接触到了阿尔伯特·爱因斯坦的狭义相对论，当时这一理论尚未被广泛接受。他在数学上的贡献帮助该理论赢得了怀疑者的认可。1914年，他与莱昂·布里渊合作研究了色散介质中电磁波的传播。他成为量子力学的奠基人之一，其贡献包括：与威尔逊共同发现索末菲尔德–威尔逊量子化规则（1915年），对玻尔原子模型的推广；引入索末菲尔德精细结构常数（1916年）；与瓦尔特·科塞尔共同发现索末菲尔德–科塞尔位移定律（1919年）；\(^\text{[41]}\)以及出版《原子结构与谱线》（Atombau und Spektrallinien，1919年），该书成为新一代理论物理学家研究原子物理与量子物理的“圣经”。\(^\text{[42]}\)该书在1919年至1924年间出版了四个版本，以纳入量子力学的最新进展，之后分为两卷出版。\(^\text{[43]}\)

1918年，索末菲尔德接替爱因斯坦出任德国物理学会主席。\(^\text{[13]}\)他的重要成就之一是创办了一份新期刊。\(^\text{[44]}\)由于DPG期刊发表的科学论文数量过于庞大，1919年学会委员会建议创立《物理学杂志》，用于刊登原创研究论文，并于1920年正式创刊。该期刊允许任何有声望的科学家在无需审稿的情况下发表文章，因此投稿与出版之间的周期极短——最快仅需两周。这极大地推动了科学理论的发展，尤其是当时德国的量子力学研究，因为这本期刊成为新一代持前卫观点的量子理论家们发表成果的首选平台。\(^\text{[45]}\)

在1922/1923年冬季学期，索末菲尔德受邀在威斯康星大学麦迪逊分校担任“卡尔·舒尔茨纪念物理学教授”讲座。\(^\text{[13]}\)1926年4月13日，他当选为曼彻斯特文学与哲学学会的荣誉会员。\(^\text{[46]}\)

1927年，索末菲尔德将费米–狄拉克统计应用于保罗·德鲁德提出的金属电子德鲁德模型。这一新理论解决了原始模型在预测热学性质时的诸多问题，并被称为德鲁德–索末菲尔德模型。

1928年至1929年，索末菲尔德进行了环球旅行，重要的访问地包括印度、\(^\text{[47]}\)中国、日本和美国。

索末菲尔德是一位杰出的理论物理学家。除了在量子理论中的不可替代的贡献之外，他还在物理学的其他领域有所建树，例如经典电磁理论。举例来说，他提出了关于导电地面上辐射赫兹偶极子的解法，这一方案在之后的岁月里产生了众多应用。他提出的索末菲尔德恒等式和索末菲尔德积分至今仍是解决此类问题的最常用方法。与此同时，作为索末菲尔德理论物理学派实力的象征，以及20世纪初理论物理崛起的标志，截至1928年，德语世界近三分之一的理论物理学正教授都是索末菲尔德的学生。\(^\text{[48]}\)
\begin{figure}[ht]
\centering
\includegraphics[width=6cm]{./figures/7b7c7be0b9c92309.png}
\caption{阿诺德·索末菲尔德，斯图加特，1935年} \label{fig_ANOSf_2}
\end{figure}
1935年4月1日，索末菲尔德获得了名誉教授身份。在继任者遴选过程中，他暂时担任自己的接替者，这一过程一直拖延到1939年12月1日才结束。之所以如此漫长，是因为慕尼黑大学教授团的推荐人选与德意志教育部及“德意志物理学”支持者之间存在学术与政治分歧。\(^\text{[49][50]}\)“德意志物理学”具有反犹倾向，并对理论物理，尤其是量子力学，抱有强烈的偏见。最终任命的威廉·米勒既不是理论物理学家，从未在物理学期刊发表过论文，也不是德国物理学会的成员，\(^\text{[51]}\)他接替索末菲尔德的任命被认为是对物理教育的一种嘲弄，严重不利于培养新一代物理学家。凯撒·威廉流体研究所所长路德维希·普朗特尔，以及通用电气公司研究部门主任兼德国物理学会主席卡尔·拉姆绍尔，都在写给帝国官员的信中提及此事。在普朗特尔1941年4月28日写给帝国元帅赫尔曼·戈林的信件附件中，他将这一任命称为对必要的理论物理教育的“破坏”。\(^\text{[52]}\)在拉姆绍尔1942年1月20日写给帝国部长伯恩哈德·鲁斯特的信件附件中，拉姆绍尔则断言这一任命等同于对“慕尼黑理论物理传统的毁灭”。\(^\text{[53][54]}\)

至于索末菲尔德曾经的爱国主义观点，他在希特勒上台后不久写信给爱因斯坦说道：“我可以向你保证，我们统治者对‘民族’一词的滥用，彻底让我戒掉了那种在我身上曾经如此强烈的民族情感。我现在宁愿看到德国作为一个强权消失，并融入一个和平的欧洲。”\(^\text{[55]}\)

索末菲尔德一生中获得了许多荣誉，例如：洛伦兹奖章、马克斯·普朗克奖章、厄斯特奖章，\(^\text{[56][57]}\)当选伦敦皇家学会、美国国家科学院、苏联科学院、印度科学院以及柏林、慕尼黑、哥廷根和维也纳等科学院院士，并获得罗斯托克、亚琛、加尔各答和雅典等大学授予的多项名誉学位。\(^\text{[8]}\)他于1950年被选为光学学会荣誉会员。\(^\text{[58]}\)他曾被提名诺贝尔奖84次，比任何其他物理学家都多（包括被提名82次的奥托·斯特恩\(^\text{[59]}\)），但他始终未能获奖。\(^\text{[60][61]}\)
\begin{figure}[ht]
\centering
\includegraphics[width=6cm]{./figures/d1b8061dc47a9103.png}
\caption{慕尼黑诺德公墓中的索末菲尔德之墓} \label{fig_ANOSf_3}
\end{figure}
1951年4月26日，索末菲尔德在慕尼黑去世，死因是与孙辈散步时遭遇车祸所受的伤害。事故发生在他位于杜南特大街6号住所附近的迪特林登大街与比德施泰纳大街拐角处。\(^\text{[62][63]}\)他被安葬在诺德公墓，距离当时的住所不远。\(^\text{[63]}\)

2004年，慕尼黑大学的理论物理中心以他的名字命名。\(^\text{[64]}\)
\subsection{著作}
\subsubsection{论文}
\begin{itemize}
\item 阿诺德·索末菲尔德，《衍射的数学理论》，Mathematische Annalen*47(2–3)，第317–374页。（1896年）。doi:10.1007/bf01447273。
\item   由Raymond J. Nagem、Mario Zampolli和Guido Sandri翻译为《衍射的数学理论》（Mathematical Theory of Diffraction*，Birkhäuser Boston，2003年），ISBN 0-8176-3604-8。
\item 阿诺德·索末菲尔德，《无线电报中波的传播》，Annalen der Physik [第4辑] 28, 665 (1909)；62, 95 (1920)；81, 1135 (1926)。
\item 阿诺德·索末菲尔德，《我教学生涯的一些回忆》，American Journal of Physics 第17卷，第5期，第315–316页（1949年）。这是他在接受1948年厄斯特奖章时的演讲。
\end{itemize}
\subsubsection{书籍}
\begin{itemize}
\item 阿诺德·索末菲尔德，《原子结构与谱线》，不伦瑞克：Friedrich Vieweg und Sohn，1919年。
\item 阿诺德·索末菲尔德，《原子结构与谱线》，由Henry L. Brose根据德文第三版翻译，伦敦：Methuen，1923年。
\item 阿诺德·索末菲尔德，《原子物理三讲》，伦敦：Methuen，1926年。
\item 阿诺德·索末菲尔德，《原子结构与谱线：波动力学补编》，不伦瑞克：Vieweg，1929年。
\item 阿诺德·索末菲尔德，《波动力学：〈原子结构与谱线〉补编》，Henry L. Brose 翻译，纽约：Dutton，1929年。
\item 阿诺德·索末菲尔德，《加尔各答大学讲授的波动力学讲义》，加尔各答大学，1929年。
\item 阿诺德·索末菲尔德、汉斯·贝特，《金属的电子理论》，载于 H. Geiger 与 K. Scheel 主编的《物理学手册》，第24卷第2分册，第333–622页（柏林：Springer，1933年）。这篇近300页的章节后来作为单行本出版：《金属的电子理论》（Springer，1967年）。
\item 阿诺德·索末菲尔德，《力学——理论物理讲义·第一卷》，莱比锡：Akademische Verlagsgesellschaft Becker & Erler，1943年。
\item 阿诺德·索末菲尔德，《力学——理论物理讲义·第一卷》，由 Martin O. Stern 根据德文第四版翻译，纽约：Academic Press，1964年。
\item 阿诺德·索末菲尔德，《可变形介质力学——理论物理讲义·第二卷》，莱比锡：Akademische Verlagsgesellschaft Becker & Erler，1945年。
\item 阿诺德·索末菲尔德，《可变形体力学——理论物理讲义·第二卷》，由 G. Kuerti 根据德文第二版翻译，纽约：Academic Press，1964年。
\item 阿诺德·索末菲尔德，《电动力学——理论物理讲义·第三卷》，Klemm Verlag，1948年。
\item 阿诺德·索末菲尔德，《电动力学——理论物理讲义·第三卷》，Edward G. Ramberg 翻译，纽约：Academic Press，1964年。
\item 阿诺德·索末菲尔德，《光学——理论物理讲义·第四卷》，Dieterich’sche Verlagsbuchhandlung，1950年。
\item 阿诺德·索末菲尔德，《光学——理论物理讲义·第四卷》，由 Otto Laporte 与 Peter A. Moldauer 根据德文第一版翻译，纽约：Academic Press，1964年。
\item 阿诺德·索末菲尔德，《热力学与统计学——理论物理讲义·第五卷》，由弗里茨·博普（Fritz Bopp）与约瑟夫·迈克斯纳编辑，Diederich’sche Verlagsbuchhandlung，1952年。
\item 阿诺德·索末菲尔德，《热力学与统计力学——理论物理讲义·第五卷》，由 F. Bopp 与 J. Meixner 编辑，J. Kestin 翻译，纽约：Academic Press，1964年。
\item 阿诺德·索末菲尔德，《物理中的偏微分方程——理论物理讲义·第六卷》，Dieterich’sche Verlagsbuchhandlung，1947年。
\item 阿诺德·索末菲尔德，《物理中的偏微分方程——理论物理讲义·第六卷》，Ernest G. Straus 翻译，纽约：Academic Press，初版1949年，第二版1953年；同时为 AP *Pure and Applied Mathematics* 系列第1号。
\item 费利克斯·克莱因、阿诺德·索末菲尔德，《陀螺理论》[四卷本]，莱比锡：Teubner，1897年。
\item 费利克斯·克莱因、阿诺德·索末菲尔德，《陀螺理论·第一卷》，Raymond J. Nagem 与 Guido Sandri 翻译，波士顿：Birkhäuser，2008年。ISBN 0817647201。
\end{itemize}
\subsection{另见}
\begin{itemize}
\item 以阿诺德·索末菲尔德命名的事物列表
\end{itemize}
\subsection{参考文献}
\begin{enumerate}
\item 马克斯·玻恩 (1952)。《阿诺德·约翰内斯·威廉·索末菲尔德，1868–1951》。英国皇家学会会员讣告 ，第8卷(21)：274–296。doi:10.1098/rsbm.1952.0018。JSTOR 768813。S2CID 161998194。
\item 《Arnold Sommerfeld | Quantum Mechanics, Atomic Theory, Mathematics | Britannica》。[www.britannica.com。2024年4月22日。2024年5月12日检索。](http://www.britannica.com。2024年4月22日。2024年5月12日检索。)
\item  O'Connor, J.J.; Robertson, E.F. 《Arnold Sommerfeld (1868–1951) – Biography – MacTutor 数学史档案》。2022年10月31日存档。2022年10月31日检索。
\item 《Sommerfeld, Arnold (Johannes Wilhelm)》。科学传记全书 。Encyclopedia.com。2022年10月31日存档。2022年10月31日检索。
\item Michael Eckert (2013)。《阿诺德·索末菲尔德：科学、人生与动荡的时代 1868–1951》。Tom Artin 翻译。纽约：Springer。ISBN 9781461474609。OCLC 832272564。
\item 数学家谱系项目（The Mathematics Genealogy Project，记录存档于2006年10月17日 Wayback Machine）列出费迪南德·冯·林德曼为索末菲尔德的博士论文导师。而卡西迪（Cassidy, 1992, 第100–101页）则认为保罗·福尔克曼是索末菲尔德的导师，并提供了参考文献。其他一些作者则根据索末菲尔德的学术能力，提供了可以在二者之间进行判断的信息。索末菲尔德《教授资格论文》的英文译本（Arnold Sommerfeld，由Raymond J. Nagem、Mario Zampolli 和 Guido Sandri 翻译，《衍射的数学理论》，Birkhäuser Boston，2003，第1–2页）显示，索末菲尔德在博士论文中引用了柯尼斯堡大学的14位教师，并对他们全部表示感谢，但在致谢语中特别点名了林德曼。荣克尼克尔（Jungnickel, 1990b, 第144–148页和第157–160页）对与福尔克曼相关的诸多问题有揭示性论述。他本人研究不多，对物理学家吸引力有限，发表的论文寥寥无几，作为物理教师更像是一名“科普者”。虽然索末菲尔德在柯尼斯堡确实上过福尔克曼理论物理研究所的课程，但他依靠的却是福尔克曼的助手埃米尔·维歇特，而非福尔克曼本人。索末菲尔德与维歇特关系密切，后者给了他许多重要的影响。Wilfried Schroeder 曾发表过索末菲尔德与维歇特之间的早期书信（Arch. hist. ex. sci., 1984）。在19世纪末至20世纪初，理论物理的正教授职位在德国只有四个：柯尼斯堡（福尔克曼）、哥廷根、柏林（马克斯·普朗克），以及慕尼黑——该职位自路德维希·玻尔兹曼于1894年离开后一直空缺，直到1906年索末菲尔德被任命时才得到补缺。在1899年关于理论物理学地位的评论中，福格特只提到了普朗克、威廉·维恩、保罗·德鲁德以及索末菲尔德。而维恩在1898年写给索末菲尔德的信中，其评价与福格特类似：他只提到了柏林和哥廷根的教授职位。考虑到慕尼黑职位当时尚未填补，却完全没有提及柯尼斯堡福尔克曼的教授职位，这一“遗漏”极为显眼，且耐人寻味。
\item Mehra，第1卷，第1部分，1982年，第106页。
\item 《索末菲尔德传记》，2007年9月29日存档于 Wayback Machine —— MacTutor 数学史档案。
\item 《阿诺德·索末菲尔德》，2006年10月17日存档于 Wayback Machine —— 数学家谱系项目。索末菲尔德博士论文题目：《数理物理中的任意函数》。
\item 《阿诺德·索末菲尔德传记》，2006年9月27日存档于 Wayback Machine —— 美国哲学学会。
\item 阿诺德·索末菲尔德，《衍射的数学理论》，Raymond J. Nagem、Mario Zampolli、Guido Sandri 翻译，Birkhäuser Boston，2003年，ISBN 0-8176-3604-8。
\item 索末菲尔德教授资格论文题目：《衍射的数学理论》。
\item 《索末菲尔德项目》，2012年8月5日存档于 archive.today —— 莱布尼茨科学计算中心。
\item 阿诺德·索末菲尔德，《力学——理论物理讲义·第一卷》，由 Martin O. Stern 根据德文第四版翻译，Academic Press，1964年，第v–x页。（含保罗·彼得·埃瓦尔德序言及索末菲尔德前言。）
\item 馆长是大学中常驻的政府代表。
\item 邓肯·道森，1979年10月1日，《摩擦学人物：威廉·贝特·哈代（1864–1934）与阿诺德·索末菲尔德（1868–1951）》，润滑技术杂志，第101卷（4）：393–397。doi:10.1115/1.3453381。ISSN 0022-2305。
\item 荣克尼克尔，1990b，第274页、第277–278页和第281–285页。
\item 荣克尼克尔，1990b，第157–160页，第254页起，第304页起，以及第384页起。
\item 卡西迪，1992，第104页。
\item 索末菲尔德的学生可以按照类型分类，即在他指导下的学习路径。（部分学生和博士后的相关脚注请参见正文。）
\begin{itemize}
\item 博士生：卡尔·阿普费尔巴赫、赫尔曼·布吕克、保罗·索福斯·爱普斯坦、约翰内斯·菲舍尔、瓦尔特·弗兰茨、赫尔伯特·弗勒利希、欧文·福厄斯、卡尔·格利彻、弗雷德里克·格罗弗、恩斯特·吉耶尔曼、维尔纳·海森堡、德米特里乌斯·洪德罗斯、赫尔穆特·黑内尔、路德维希·霍普夫、阿尔弗雷德·朗德、赫尔伯特·朗、奥托·拉波尔特、卡尔·迈斯纳、约瑟夫·迈克斯纳、海因里希·奥特、沃尔夫冈·泡利、爱德华·兰伯格、瓦伦丁·沙伊德尔、奥托·舍尔策、鲁道夫·泽利格尔、恩斯特·施蒂克尔伯格以及阿尔布雷希特·云索尔德。
\item 博士学位在其他地方完成要求者：莱昂·布里渊、尤金·芬伯格、鲁道夫·佩耶尔斯以及弗朗西斯·G·斯莱克。
\item 博士学位由索末菲尔德担任副导师（主要导师括号内注明）：弗里德里希·布尔迈斯特（Friedrich Burmeister – 导师：Hugo von Seeliger）、瓦尔特·海特勒（Walter Heitler – 导师：Karl Herzfeld）、卡尔·赫茨菲尔德（Karl Herzfeld – 导师：Friedrich Hasenöhrl）、赫尔曼·马奇（Herman March – 导师：Wilhelm Röntgen）、库尔特·乌尔班（Kurt Urban – 导师：Wilhelm Rabe）、卡尔·泽巴赫（Karl Seebach – 导师：Heinrich Tietze）以及布鲁诺·蒂林（Bruno Thüring – 导师：Alexander Wilkens）。
\item 博士与教授资格：卡尔·贝歇特、汉斯·贝特、彼得·德拜、保罗·彼得·埃瓦尔德、阿道夫·克拉策尔、威廉·伦茨、路德维希·瓦尔德曼、海因里希·韦尔克以及格雷戈尔·文策尔。
\item 教授资格：瓦尔特·科塞尔、马克斯·冯·劳厄以及沃伊切赫·鲁比诺维奇。
\item 博士后：威廉·阿利斯、博伊德·巴特利特、理查德·贝克尔、爱德华·康登、卡尔·埃卡特、威廉·V·休斯顿、埃德温·C·肯布尔、菲利普·M·莫尔斯、卡雷尔·尼森、莱纳斯·鲍林、伊西多·拉比、霍华德·罗伯逊以及弗里茨·绍特。
\item 亚琛本科生：瓦尔特·罗戈夫斯基。
\end{itemize}
\item 阿诺德·索末菲尔德的学生，2006年10月17日存档于 Wayback Machine —— 数学家谱系项目；以及《阿诺德·索末菲尔德——交流与学校教育》，2007年6月15日存档于 Wayback Machine。
\item 鲁道夫·佩耶尔斯于1926年至1928年在索末菲尔德门下攻读博士，随后转至莱比锡大学，在沃尔夫冈·泡利指导下完成博士学位，并于1929年获授。参见：美国哲学学会作者目录：Peierls，2007年2月5日存档于 Wayback Machine。
\item 尤金·芬伯格也曾在索末菲尔德门下攻读博士，之后于1933年在哈佛大学埃德温·C·肯布尔的指导下完成博士学位。
\item 卡尔·迈斯纳在慕尼黑随索末菲尔德学习一年后，返回蒂宾根，以便在弗里德里希·帕申的指导下研究光谱学，并于1915年获得博士学位。参见：K. W. Meissner 书评：《阿诺德·索末菲尔德〈光学——理论物理讲义·第四卷〉》（由 Otto Laporte 与 Peter A. Moldauer 根据德文第一版翻译），美国物理学杂志23卷(7)，477–478（1955年）。作者在文中说明，他曾于1912年听过索末菲尔德的讲座，尤其是关于光学的课程。
\item 1912年至1913年间，莱昂·布里渊曾在索末菲尔德门下攻读研究生。随后，他于1920年在巴黎大学，在保罗·朗之万的指导下获得了国家博士学位。参见：美国哲学学会作者目录：Brillouin，2007年2月5日存档于 Wayback Machine。
\item 保罗·鲍林通过1925–1926年的美国国家研究委员会奖学金，以及1926–1927年的古根海姆基金会奖学金，先后在索末菲尔德、苏黎世的埃尔温·薛定谔以及哥本哈根的尼尔斯·玻尔门下从事博士后研究。参见：诺贝尔奖传记：Pauling，2007年6月20日存档于 Wayback Machine。另见：阿诺德·索末菲尔德，《我教学生涯的一些回忆》，美国物理学杂志第17卷，第315–316页（1949年）。在这篇文章中，索末菲尔德特别提及其（博士后）学生包括美国人莱纳斯·鲍林、爱德华·U·康登和伊西多·拉比。
\item 1927年获得博士学位后，伊西多·拉比在奖学金的资助下前往欧洲两年，先后在索末菲尔德、尼尔斯·玻尔、沃尔夫冈·泡利、奥托·斯特恩以及维尔纳·海森堡门下从事博士后研究。参见：Isidor Isaac Rabi – Biographical，2012年10月2日存档于 Wayback Machine。另见：阿诺德·索末菲尔德，《我教学生涯的一些回忆》，美国物理学杂志第17卷，第315–316页（1949年）。在这篇文章中，索末菲尔德特别提到，他的（博士后）学生中包括美国人莱纳斯·鲍林、爱德华·U·康登和伊西多·拉比。另见：I. I. Rabi，《量子力学早期的故事》，由R. Fraser Code 翻译与编辑，Physics Today 第8期，第36–41页（2006年）。在文中，拉比回忆了他作为索末菲尔德博士后的经历。
\item Walker, 1995，第73页。冯·劳厄于1906年完成教授资格。
\item 在1930–1931学年，威廉·阿利斯前半学期随索末菲尔德学习，后半学期则在剑桥大学。他与菲利普·M·莫尔斯一同旅行。参见：Morse 1978，第100页。
\item 在获得博士学位后，爱德华·康登于1926年和1927年，凭借洛克菲勒基金会资助的国家研究委员会奖学金，在慕尼黑随索末菲尔德，以及在哥廷根随马克斯·玻恩从事博士后研究。参见：美国物理学会：Edward Condon，2006年12月6日存档于 Wayback Machine；以及美国哲学学会 – MOLE: Condon，2006年9月27日存档于 Wayback Machine。另见：阿诺德·索末菲尔德，《我教学生涯的一些回忆》，美国物理学杂志第17卷，第315–316页（1949年）。在文中，索末菲尔德特别提及他的（博士后）学生包括美国人莱纳斯·鲍林、爱德华·U·康登和伊西多·拉比。
\item 1927年至1928年，卡尔·埃卡特获得古根海姆基金会奖学金，并利用该资助前往德国，先后在慕尼黑的路德维希–马克西米利安大学随阿诺德·索末菲尔德，以及在莱比锡大学随维尔纳·海森堡从事博士后研究。参见：Eckart 传记 – 美国国家科学院出版社；以及作者目录：Eckart，2007年2月5日存档于 Wayback Machine —— 美国哲学学会。另见：阿诺德·索末菲尔德，《我教学生涯的一些回忆》，美国物理学杂志，第17卷（第5期），315–316页（1949年）。
\item 埃德温·C·肯布尔于1927–1928年间前往慕尼黑和哥廷根，分别随索末菲尔德和马克斯·玻恩学习并开展研究。
\item 1927年至1928年，威廉·V·休斯顿获得古根海姆基金会奖学金，并利用该资助前往德国，先后在慕尼黑路德维希–马克西米利安大学随索末菲尔德，以及在莱比锡大学随维尔纳·海森堡从事博士后研究。参见：Houston 传记，2007年3月24日存档于 Wayback Machine —— 美国国家科学院出版社。另见：阿诺德·索末菲尔德，《我教学生涯的一些回忆》，美国物理学杂志第17卷（第5期），315–316页（1949年）。
\item 卡尔·费迪南德·赫茨菲尔德曾在慕尼黑大学与索末菲尔德及卡齐米日·法扬斯共事，起初担任私人讲师（Privatdozent，1919–1925年），随后担任特别教授（extraordinarius professor，1925–1926年）。参见：Karl Ferdinand Herzfeld 1892–1978，Joseph F. Mulligan 撰写的传记回忆录，美国国家科学院出版社，2001年，2015年9月12日存档于 Wayback Machine。
\item 保罗·柯克帕特里克，作为奖项委员会主席，在《美国物理学杂志》17卷（第5期）312–314页（1949年）中发表推荐致辞。在这篇文章中，提到了以下索末菲尔德的学生：威廉·V·休斯顿、卡尔·贝歇特、奥托·舍尔策、奥托·拉波尔特、莱纳斯·鲍林、卡尔·埃卡特、格雷戈尔·文策尔、彼得·德拜以及菲利普·M·莫尔斯。
\item 莫尔斯，菲利普·M.（1978）。《在开端之中：一位物理学家的生活》，第二次印刷版。麻省理工学院出版社，第100页。ISBN 0262131242。
\item I. I. 拉比，由 R. Fraser Code 翻译与编辑，《量子力学早期的故事》，Physics Today，（第8期），第36–41页（2006），第38页。
\item 鲁比诺维奇，于1916年至1918年在慕尼黑学习。
\item 荣克尼克尔，1990b，第284页，引自该页脚注100中的参考文献。
\item 荣克尼克尔，1990b，第283页。
\item 梅赫拉，1982年，第1卷，第1部分，第330页。
\item 克拉格，2002年，第155页。
\item 埃克特，迈克尔（2013年9月10日），《索末菲尔德的〈原子结构与谱线〉》，载于 研究与教学：通过教材看量子物理史，MPRL – Studies，柏林：马克斯·普朗克学会，ISBN 978-3-945561-24-9，2024年4月17日检索。
\item 卡西迪，1992年，第106页。
\item 克拉格，2002年，第150页与第168页。
\item 《曼彻斯特回忆录与会议录》。
\item 拉金德·辛格，〈阿诺德·索末菲尔德——德国印度物理学的支持者〉，Current Science 第81卷第11期，2001年12月10日，第1489–1494页，2007年3月27日存档于 Wayback Machine。
\item 拜尔亨，1977年，第9页，引自以下文献：马克斯·玻恩，《作为学派奠基人的索末菲尔德》，Die Naturwissenschaften第16卷，第1036页（1928年）。
\item 拜尔亨，1977年，第150–167页。
\item 1935年，慕尼黑大学教授团拟定了一份候选人名单，用以接替索末菲尔德担任理论物理学正教授及理论物理研究所所长。名单上有三人：维尔纳·海森堡（Werner Heisenberg，1932年诺贝尔物理学奖得主）、彼得·德拜（Peter Debye，1936年诺贝尔化学奖得主）以及理查德·贝克尔——他们都是索末菲尔德的前学生。慕尼黑教授团坚决支持这些候选人，并将海森堡列为首选。\\
然而，“德意志物理学”的支持者以及德意志教育部的某些势力则有他们自己的人选名单，由此引发了一场长达四年多的斗争。在此期间，海森堡遭到了“德意志物理学”支持者的猛烈攻击，其中一次攻击发表于党卫队（SS）机关报《黑色军团报》，该报由海因里希·希姆莱主管。甚至有一次，海森堡的母亲不得不拜访希姆莱的母亲，试图帮助平息此事。两位母亲彼此相识，是因为海森堡的外祖父与希姆莱的父亲都曾担任校长，并且同为巴伐利亚登山俱乐部成员。\\
最终，这场“海森堡事件”得以解决，但结局却是政治上的胜利而非学术标准的胜利。1939年12月1日，威廉·米勒接替了索末菲尔德。\\
参见：Beyerchen, 1977, 第153–167页；Cassidy, 1992, 第383–387页；Thomas Powers，《海森堡的战争：德国原子弹的秘密历史》（Heisenberg's War: The Secret History of the German Bomb，Knopf，1993年），第40–43页；Hentschel, 1996, 第176–177页；以及Samuel A. Goudsmit，《ALSOS》（Tomash Publishers, 1986年），第117–119页。

\item 拜尔亨（Beyerchen），1977年，第166页。
\item 亨切尔（Hentschel），1996年，第265页。见：亨切尔1996年文献第85号，第261–266页。
\item 亨切尔（Hentschel），1996年，第291页。见：亨切尔1996年文献第93号，第290–292页。
\item 亨切尔（Hentschel），1996年，附录F；参见卡尔·拉姆绍尔（Carl Ramsauer）与路德维希·普朗特尔（Ludwig Prandtl）的条目。
\item 《MacTutor 数学史档案：阿诺德·索末菲尔德（1868–1951）》。2019年12月2日检索。
\item 柯克帕特里克，保罗（Paul Kirkpatrick）（1949）。《推荐致辞》，由奖项委员会主席柯克帕特里克教授发表。*美国物理学杂志*（American Journal of Physics），第17卷（第5期）：312–314。Bibcode:1949AmJPh..17..312.. doi:10.1119/1.1989584。
\item 索末菲尔德，阿诺德（Arnold Sommerfeld）（1949）。《我教学生涯的一些回忆》（*Some Reminiscences of My Teaching Career*）。*美国物理学杂志*，第17卷（第5期）：315–316。Bibcode:1949AmJPh..17..315S。doi:10.1119/1.1989585。S2CID 122912752。
\item 《Arnold Sommerfeld | Optica》。[www.optica.org。2024年9月23日检索。](http://www.optica.org。2024年9月23日检索。)
\item 《提名数据库》（Nomination Database）。[www.nobelprize.org。2016年11月6日存档。2018年5月4日检索。](http://www.nobelprize.org。2016年11月6日存档。2018年5月4日检索。)

\item 《提名数据库》（*Nomination Database*）。[www.nobelprize.org。2018年2月12日存档。2018年5月4日检索。](http://www.nobelprize.org。2018年2月12日存档。2018年5月4日检索。)
\item 克劳福德，伊丽莎白（Elisabeth Crawford）（2001年11月）。〈1901–1950年诺贝尔群体：科学精英的剖析〉（*Nobel population 1901–50: anatomy of a scientific elite*）。*Physics World*。2006年2月3日存档。
\item 《德克萨斯大学：塞西尔·德维特–莫雷特（Cécile DeWitt-Morette，1922年12月21日－2017年5月8日）》，*德克萨斯大学奥斯汀分校物理系历史*。信件作者：沃尔夫冈·泡利（Prof. Dr. W. Pauli）。2019年6月20日存档。2019年2月19日检索。信中写道：“我刚刚经历了一件令人悲伤的事：我的老导师索末菲尔德本周在慕尼黑因交通事故去世，享年82岁。”
\item 张秀中（Sooyojng Chang）（2005）。《物理学家的学术谱系》（*Academic Genealogy of Physicists*）。首尔国立大学出版社，第121页。
\item 《阿诺德·索末菲尔德理论物理中心 – 首页》（*Arnold Sommerfeld Center for Theoretical Physics – Home*）。2007年2月10日存档。2007年2月12日检索。——阿诺德·索末菲尔德理论物理中心。
\end{enumerate}
\subsection{延伸阅读}
\begin{itemize}
\item 本茨，乌尔里希，《阿诺德·索末菲尔德：原子时代门槛上的教师与研究者 1868–1951》，Wissenschaftliche Verlagsgesellschaft，1975年。
\item 拜尔亨，艾伦·D.，《希特勒治下的科学家：第三帝国的政治与物理学界》，耶鲁大学出版社，1977年。
\item 玻恩，马克斯，《阿诺德·约翰内斯·威廉·索末菲尔德，1868–1951》，英国皇家学会会员讣告，第8卷，第21期，第274–296页（1952年）。
\item 卡西迪，大卫·C.，《不确定性：维尔纳·海森堡的生活与科学》，W. H. Freeman and Company，1992年，ISBN 0-7167-2503-7。（由于维尔纳·海森堡是索末菲尔德的博士生，本书虽为间接资料，但关于索末菲尔德的内容相当丰富且有详实记录。
\item 埃克特，迈克尔，《阿诺德·索末菲尔德：原子物理学家与文化使者 1868–1951。传记》，德国博物馆，Wallstein Verlag，2013年。
\item 埃克特，迈克尔，汤姆·阿廷译，《阿诺德·索末菲尔德：科学、人生与动荡时代，1868–1951》，Springer，2013年。
\item 埃克特，迈克尔（Michael Eckert），〈科学中的宣传：索末菲尔德与金属电子理论的传播〉，历史研究：物理与生物科学，第17卷，第2期，第191–233页（1987年）。
\item 埃克特，迈克尔，〈数学、实验与理论物理：索末菲尔德学派的早期岁月〉，视角中的物理学，第1卷，第3期，第238–252页（1999年）。
\item 亨切尔，克劳斯，安·M·亨切尔，《物理学与国家社会主义：一手资料选集》，Birkhäuser，1996年。
\item 荣克尼克尔，克里斯塔与拉塞尔·麦考马克，《自然的理智掌控：从欧姆到爱因斯坦的理论物理学》，第1卷：《数学的火炬，1800–1870》，芝加哥大学出版社，平装版，1990a。ISBN 0-226-41582-1。
\item 荣克尼克尔，克里斯塔与拉塞尔·麦考马克，《自然的理智掌控：从欧姆到爱因斯坦的理论物理学》第2卷：《如今强大的理论物理学，1870–1925》，芝加哥大学出版社，平装版，1990b。ISBN 0-226-41585-6。
\item 康特，霍斯特，〈阿诺德·索末菲尔德——交流与学校教育〉，见于Fuchs-Kittowski, Klaus; Laitko, Hubert; Parthey, Heinrich; Umstätter, Walther（编），《科学与数字图书馆：科学研究年鉴 1998》，第135–152页，科学研究学会出版社，2000年。
\item 柯克帕特里克，保罗，〈奖项委员会主席柯克帕特里克教授的推荐致辞〉，美国物理学杂志），第17卷，第5期，第312–314页（1949年）。这是在1949年1月28日为阿诺德·索末菲尔德颁发1948年厄斯特奖章（以表彰其对物理教学的杰出贡献）前的致辞。
\item 克拉格，赫尔格，《量子世代：二十世纪物理学史》，普林斯顿大学出版社，第五次印刷及平装版首次印刷，2002年，ISBN 0-691-01206-7。
\item 库恩，托马斯·S.、约翰·L·海尔布隆、保罗·福尔曼、利尼·艾伦，《量子物理史资料》，美国哲学学会，1967年。
\item 梅赫拉，贾格迪什与赫尔穆特·雷兴贝格（，《量子理论的历史发展》第一卷第一部分：《普朗克、爱因斯坦、玻尔与索末菲尔德的量子理论 1900–1925：它的奠基与困境的出现》，Springer，1982年，ISBN 0-387-95174-1。
\item 鲍林，莱纳斯，〈阿诺德·索末菲尔德：1868–1951〉，科学，第114卷，第2963期，第383–384页（1951年）。
\item 辛格，拉金德，〈阿诺德·索末菲尔德——德国印度物理学的支持者〉，Current Science，第81卷第11期，2001年12月10日，第1489–1494页，2007年3月27日存档于 Wayback Machine。
\item 沃克，马克，《纳粹科学：神话、真相与德国原子弹》，Persius出版社，1995年，ISBN 0-306-44941-2。
\end{itemize}